% This file was converted to LaTeX by Writer2LaTeX ver. 1.0.2
% see http://writer2latex.sourceforge.net for more info
\documentclass[a4paper,twocolumn]{article}
\usepackage{lastpage}
\usepackage[utf8]{inputenc}
\usepackage[T1]{fontenc}
\usepackage[english]{babel}
\usepackage{times}
\usepackage{amsmath}
\usepackage{amssymb,amsfonts,textcomp}
\usepackage{color}
\usepackage[scale={.85,.9}]{geometry}
\usepackage{array}
\usepackage{mdwlist}
\usepackage{hhline}
\usepackage{hyperref}
\hypersetup{pdftex, colorlinks=true, linkcolor=blue, citecolor=blue, filecolor=blue, urlcolor=blue, pdftitle=War on Fraud, pdfauthor=Pieter Van den Hombergh, pdfsubject=Avoid fraud on Fontys Campus Venlo, pdfkeywords=Fraud computers electronic devices examination rules}
\usepackage[pdftex]{graphicx}
% Footnote rule
\setlength{\skip\footins}{0.119cm}
\renewcommand\footnoterule{\vspace*{-0.018cm}\setlength\leftskip{0pt}\setlength\rightskip{0pt plus 1fil}\noindent\textcolor{black}{\rule{0.25\columnwidth}{0.018cm}}\vspace*{0.101cm}}
% Pages styles
\makeatletter
\newcommand\ps@Standard{
  \renewcommand\@oddhead{}
  \renewcommand\@evenhead{}
  \renewcommand\@oddfoot{\selectlanguage{english} \thepage{}/\pageref{LastPage}}
  \renewcommand\@evenfoot{\@oddfoot}
  \renewcommand\thepage{\arabic{page}}
}
\makeatother
\pagestyle{Standard}
\title{War on Fraud}
\author{Pieter Van den Hombergh}
\date{2014-03-09}
\title{War on fraud}
\begin{document}
\maketitle

\abstract{Some thoughts on fraud, possibilities of fraud and countermeasures.}

\textbf{Teaching is fun, examination is not.}

\textit{However, examinations cannot be avoided because the
  credibility of a degree would be impacted  when  not
  examining or by examining against  too  low standards
  or in a not fraud resistant setting.}


Looking from the (honest) candidate{\textquotesingle}s point of view:


The examiner/teacher should always consider that the exam setting for
the student is a stress-full environment, because a lot can depend on
the outcome, even when the preparation to the exam ideally took the
whole semester of well spent study time. A difficulty ratio D is
comparing the difficulty X of the tasks in the exam setting versus the
difficulty P during the exercise or practical setting.  Ideally D = X/P
should be about 80\%. Exam tasks with an average difficulty level
above 100\% could be the cause of low study efficiency\footnote{NL:
  rendement} and are at the same time unfair. This implies that
introducing new concepts, like new mathematical concepts in an exam
should be declared forbidden.

\section*{Fairness: Essential for every
  examination}

Written exams are the most common form of examination because of the
affordable ratio of number of tests taken per teacher/examiner per unit
of time. Because all the candidates ideally take the same exam at the
same time\footnote{Taking the same exam at different times invariably
  leads to leaking information, making the result for the candidates
  tested at a later time biased towards a higher grade. If resource
  capacity dictates several examination times for the same subject, the
  best one can do is to sort the students by talent and let the lesser
  students go first, because the change that they will  leak useful
  information is less likely.} many candidates can be tested in one time
slot.

An examination without proper supervision is no examination. Period.

The invigilator or supervisor during an examination is responsible
minimising the possibility of fraud. In the current culture, the
candidates are effectively allowed to take all kinds of unneeded
attributes to the examination settings. Such privileges are often
explained with all kinds of plausibilities.

All a candidate requires during a written examination is:

\begin{enumerate*}
\item A set of tasks/questions on paper.
\item A set of answer sheets or examination papers.
\item A pen on pencil to write.
\item An eraser.
\item Some drawing aids (ruler, compass, geo triangle) are optional for those
  examinations in which drawings should be made.
\end{enumerate*}

Nothing else is needed.

\subsection*{Scenarios/ kinds of fraud}
In all following scenarios the candidates have names starting with the
letter A, B or C as is common in the field of information technology.
In the scenarios where two (or more) candidates are involved, candidate
A is the one requesting help and B or C are the collaborators in crime.
T is the teacher, Q is the invigilator. The Oracle is the source of all
good answers, being the study books, a stolen and correctly made
examination or simply Google and friends with someone feeding it with
the proper questions.

\subsubsection*{The Arsenal of  Alice and Bob and co.}

The tools that the candidates have to their disposal for their actions
in crime are:

\begin{enumerate*}
\item Passive communication. Alice reads the answer sheet of Bob without Bob
  knowing or will fully allowing. Bob could be considered naive but does
  not necessarily collaborate.
\item Active communication:

  \begin{enumerate*}
  \item \textbf{From outside the examination room.}
    Alexia and Carl have set up an agreement. Carl has access to the Oracle
    and communicates the correct answers to Alexia. The communication
    channels could be placement of signs, for instance the placements of
    shawls on the windowsill of papers on a bulletin board. Modern
    communication includes a blue-tooth in ear headset connected to
    Alexia{\textquotesingle}s mobile phone in her purse. Even more
    elaborate and expensive could be something like \textbf{google glass},
    which would allow Carl to read the questions at the same pace as Alexia
    does and telling the correct choice through the google glass bone
    speaker, which cannot be heard by out-standers.
  \item \textbf{Inside examination room.} Both Albert and Barny take the same
    exam at the same time in the same room. The exam contains a set of
    multiple choice questions in the same question and choice order (copier
    exam papers). Barny sits in a position relative to Albert, such that
    Albert can observe Barny without Barny or Albert making conspicuous
    movements. The communication they use is sending question number and
    correct choice by placing attributes on his desk in certain positions.
    This can be combined with signalling with finger poses. With 5 fingers
    one can communicate 32 different codes, assuming sufficient dexterity.
    This kind of communication is very common in base ball, in particular
    in the static game situations such as before a pitch. Of course any
    kind of pose of the body or body members can be used and interpreted as
    a signal. For instance supporting the head with the left arm means a,
    with the right arm b etc. Using modern technology (bluetooth in ear
    headset, mobile phone) would be a bit hard on Barny{\textquotesingle}s
    side, who has the least to gain, because in most cases Barny would have
    to talk, type or write in another place than on the allowed papers. If
    however Albert can interpret Barny{\textquotesingle}s sighs and puffs ,
    a new channel is open. Having something like Google Glass on both side
    could be very luxurious indeed. Then tapping once head with a
    (simplified) morse code is possible too.
    \begin{figure}
      \centering
        \includegraphics[width=.7\linewidth]{figures/handshowing27.pdf}
      \caption[Hand signalling the number 27.]{Hand signalling the
        number 27. Informatics students can help with the decoding of such sign language.
}
    \end{figure}
  \item \textbf{The temporal axis: }Across exam settings. Some attributes, like
    dictionaries, are provided to the candidates but are institute owned.
    One should consider that such objects can be used to carry information
    to a next examination (of the same type). Regular inspection might help
    to detect such information leaks. Take for instance the underlining of
    words or making ear marks in the pages.
  \end{enumerate*}
\end{enumerate*}

The above scenarios apply in normal pencil/pen and paper exams. Things
that can be used in these scenarios are pencil (placement of pencil
case, objects in and out of the pencil case), drinking bottles,
placement on the table, progress as in question number indicated by the
fluid level in the bottle. Let your imagination go wild and you can
come up with multiple other communication protocols.

As to drinking bottles: if allowed at all, they should contain
uncoloured water in a fully transparent bottle \textbf{without}\ label.
Otherwise the label can contain illegal information even on the label
glue side. Text on the glue side can be read through the bottle with
the luxury of magnification.

In case the candidates are allowed electronic equipment like
programmable calculators or personal computers or tablets, the number
of ways to communicate explodes. In addition, these can be combined
with that which has been proven to work in a normal written exam, as
mentioned in the previous paragraphs. If the electronic devices are
owned by the Anton and Björn, they can be prepared. In the case that
they themselves lack the competence to do the preparation, some smarter
but corrupt Charles might be willing to lend a helping hand.

\begin{enumerate*}
\item Anton has  illegal helper programs or data installed. He might
  get caught if Q sees that. So Charles created a screen shot of the
  allowed programs, say Excel or N\@tschool  of a screen showing
  inconspicuous content like the N\@Tschool front page or an almost
  empty excel sheet. He also installed two keyboard short cuts. One to
  put the screen shot on top of the active screen, covering all else
  up and one to remove said cover up. Now Anton can use the short cuts
  when Q comes too near or is in a position where he might be able to
  observe any illegalities. 
\item Annie has a \textit{``multi desktop''} extension installed on her
    laptop\footnote{A standard feature in Mac OSX and most Linux desktops,
      available as plug-in in MS Windows\par }. With a keyboard short cut she
    can quickly switch between an allowed and an illegal screen. Of course
    Annie is well trained on when to swap the screens.
\item Adrian and Bauke have agreed to use a certain drop box/ google drive
    account or subversion/git repository to share information. They might
    even let more candidates in, like their whole class. Partners in crime
    revisited.
\item Bertrand and Arjen have set up an ad hoc network via WiFi or blue-tooth,
    so they can share information. They where helped by Curt, an
    informatics drop out. Since they both have a camera and a microphone
    installed, they can use facial mimics and sighs and puffs to
    communicate. Of course they also use an in ear blue tooth head set, so
    they do not need to turn up the speakers to hear what the other puffs
    or sighs.
\end{enumerate*}

I am sure you could add additional possibilities.

\section*{Countermeasures}

For written exams (the paper and pen/pencil type) the following should
apply.

\subsection*{Provide tools and materials}

\begin{enumerate*}
\item All paper, including scrap paper should be provided by the institute
  through the invigilator at the start of the examination session. All
  paper, every sheet, should be marked with a sign, coupling it to the
  examination even. This includes scrap paper, which should be marked
  with for instance a date stamp.
  
  \begin{enumerate*}
  \item Any other paper then the paper mentioned above is illegal.
  \end{enumerate*}
\item All paper, including scrap paper, questions and answer sheets, must be
  handed in, so that the exam session does not leak marked paper or
  information.
\item The writing tools (pen, eraser) and any other aid (calculator, ruler,
  compass or geo triangle) must be provided by the institute. To be
  economic, these materials can be taken in at the end. This removes the
  need of having pencil cases on the desks. Note that the invigilator
  should either have a pencil sharpener or replacement pencils at hand.
  At the SEBI team Venlo we use a red tool box to carry these materials
  to the examination room.
\item All other attributes, like pen cases, phones, bags, etc. are disallowed.
  Jackets should also be disallowed, because they typically have extra
  pockets. Best would be to prescribe the clothing as in: long trousers
  or pants (no skirts, in particular no short ones), blouse or shirt and
  sweater or jersey (depending on temperature and season).
\end{enumerate*}
\subsection*{Table settings.}

\begin{enumerate*}
\item A typical examination room has the desks arranged in row and columns, in
  which the direction of the columns is that in which candidates deeper
  in the column look at the back of  candidates in front of him. The
  distance between the left side of a desk and the right side of the desk
  at the left hand side should be at least two arms lengths (150cm) to
  minimise visual or physical access from Alan's to
  Bink's desk.
\item Ideally, two completely different examinations take place in the same
  room, same time and duration. In such setting the even desk column is
  used for candidates of course X, with say a logistics subject and the
  odd desk column is used for course Y with an informatics paper
  examination. This puts Ali and Barbie much further apart. It also
  allows lesser distance between the desk columns.
\item Use of (card board) vision shields. These are used in the SEBI
  performance assessments (see further on). Note that they may also
  impair the view of the invigilator, which must be taken into account
  when considering applicability.
\end{enumerate*}
\subsection*{Exam preparation and correction, paper exam}
\begin{enumerate*}
\item Ideally, an examination paper (tasks, answer sheet, even scrap paper) is
  prepared (printed) for the individual student, with student name and
  number, course subject and date on each sheet of paper. This is current
  practice for some SEBI exams. It also eases the handing out of the exam
  papers and verification of handing in all provided material.
\item In case of multiple choice questions, the questions and the choice
  alternatives should be randomized, in such a way that no exam has the
  same question and choice order. This mitigates most communication
  inside the room, like hand and finger placement or simply looking at
  the neighbours answer sheet. This is current practice for a number SEBI
  pencil and paper examinations and for a few since 2005.
\item To ease correction of the exam type mentioned in the earlier two points,
  a program suite can be used, which uses work flow that involves the
  creation of the exams and afterwards electronically scanning the answer
  sheets and automatically processing and scoring the examination. Such a
  software suite, auto-multiple-choice, is currently used on a number of
  exams in the SEBI Venlo team. This work flow takes quite some initial
  effort, which is largely paid back over some semesters. A bonus to this
  way of working is a build up of a collection of easily reusable
  questions per examination subject.
\end{enumerate*}
\subsection*{Instructions to invigilator}

The invigilator is there to keep the examination playing field level,
making sure that no candidate has any advantage during the exam but by
his own study preparation.

\begin{enumerate*}
\item The invigilator instructs the (individual) student where to sit. This
  makes it possible to make the table setting less predictable to the
  candidates. Where needed, the candidates can read their table number at
  the entrance of the examination room. It is also possible to place
  named table cards on the desks, but when the candidates are randomized,
  they might have to perform a linear search which may lead to a mild but
  short chaos.
\item The invigilator shall never answer any questions. In fact it is best if
  the invigilator is not knowledgeable on the examination subject so he
  cannot leak any information but stays focused on avoiding fraud. 
\item If candidates find problems with the exam or any question in the exam,
  they can make a note of that on a specific reserved area on their
  paper, the protocol area. The corrector may take these remarks into
  consideration when correcting all works.
\item \textbf{``Repairing''}\ exams on the fly during an examination, as in
  writing corrections on a white or blackboard or giving ad hoc oral
  instructions by the invigilator or the course teacher should be
  \textbf{forbidden}.
\item The candidates should bring with them and put on their desk a means of
  identification, like a student card with a photo or a passport or other
  accepted ID. This to ensure that the proper candidate is examined. This
  absolutely necessary if the candidates are unknown to the invigilator.
\end{enumerate*}
\paragraph[The twins problem]{The twins problem}
There is still a nice problem, which has many possible twists: The
identical twins problem.

Assuming that an identical twins, named say Alicia and Beyonce, take the
same course with several examinations. If the twins are easily mixed
up, it could be possible that Alicia takes the exams she is best at
with optimal preparation first turn and Beyonce does the same with the
other exams she is best at, same preparation. Both take the resit
(maybe only for some subjects) of the examinations, but now
impersonating for the other. If Alicia passed the first time for exam
X, she will most likely to also pass the resit for X, impersonating as
Beyonce during the resit, because she already knows at least the kind
of questions and maybe also some reused questions.

We currently have no other solution available then paying extra
attention in such cases.

\section*{Electronic
  examinations or performance assessment or electronic aids during
  examinations}
  In those settings where electronic aids cannot be avoided, the aids
  should be provided by the institute, or at least the institute should
  take full control of the device during the examination session. Any
  device belonging to a candidate can be prepared by the candidate or
  some accomplice. The opportunities to such fraud should be avoided at
  all cost.

\subsection*{Institute owned and prepared devices}


The best scenario would be that the institute owns and \textit{controls}
the devices that are available to the candidates during the
examination. 

\begin{enumerate*}
\item In the case of simple calculators without persistent storage, no special
  actions need to be taken.
\item In the case of programmable calculators, the institute owned device
  should be reset (maybe to factory settings) before and/or after each
  examination. 
\item In any case, the devices should be handed out at the latests possible
  moment before the examination. Otherwise (battery) covers etc. can be
  used as covert places to store information. The same applies of course
  to dictionaries provided to the students.
\end{enumerate*}
\subsection*{Computers or virtual desktops}
In case the candidates need access to applications normally running on a
personal computer (laptop), that (virtual) computer must be owned and
or \textit{controlled} by the institute, at least for the duration of
the examination. The span of control includes preparation and post
processing time, for instance to install programmes or data (pre) and
collecting candidates work (post). In case a virtual computer image is
used (think VDI or citrix client) the thin client device used to attach
to the virtual desktop instance must underlie the same rules: Being
controlled by the institute during, just before (to make sure the
controlled client is the one used) and directly after the exam (to
avoid information leakage, like taking copies of the controlling
programmes for off-line analysis).

\subsection*{Institute taking temporary control of devices, the
  SEBI-Sandbox-stick}
One way of taking control of personal computers or laptops is letting
the computers boot from a specially prepared USB-stick. Such stick
contains a complete operating system which takes full control of the
computer. Full control means control over processor(s), all of the
memory, keyboard, display, internal or attached storage and the
communication possibilities (network, wifi, blue-tooth) of the
computer. It also means that the stick must be booted on the
\textbf{bare iron} and not as some virtual machine guest OS. The SEBI
Venlo team has developed such a stick with a simple design rule:
\textit{Disable all access unless needed.} In this
case it means that the USB-stick disables read and write access to the
laptop's internal or attached storage (hard disk, ssd,
cd, other usb, memory cards). It also turns off all communication
except the wired Ethernet network. For that network, very restricted
access rules are enforced by firewall rules on the stick. The candidate
uses the USB-stick without a password or with the same password for all
sticks, to allow easy (re)production of the USB-sticks. The candidates
should be informed that the institute takes temporary control over
their device. The invigilator must ensure that the device is booted
from the USB-stick and is not rebooted during the examination session
into another mode of operation or used as a guest in a virtual machine,
whose virtualisation host can take control by the
candidate's directions.

When properly booted the candidate has access to a restricted graphical
desktop with access only to those programs that are allowed.
Installation of own programs is prohibited as is running any program
not on the PATH environment variable. The restrictions are applied by
using a restricted shell and by proper application of read and execute
rights on the programmes available. The SEBI-stick supports several
usage modes: 

\begin{enumerate*}
\item \textbf{SEBI performance assessment
    mode.} The network firewall rules allow
  access to only one exam server and those network infrastructure
  components necessary for proper network operation (name server,
  authentication server, time server, nearest router). The students pick
  up their assignment from the exam server using a version control
  protocol (subversion) and also upload their finished work to the same
  server, possibly in several versions. Using a version control system is
  common practice in software engineering. To prevent student work loss
  due to power failure, the stick provides a \textit{persistence
    layer}, meaning it has to be ``cleaned'' after each examination. The
  cleaning operation requires some equipment but on the
  other hand can be done quite fast, typically
  14 sticks in under 2 minutes, using the proper USB stick type and brand
  and preparing OS and equipment. 
\item \textbf{Browser
    only or kiosk mode.} In this scenario only
  an internet browser is used, possibly with some special
  plugin\footnote{Note that the SEBI stick is Linux based, so internet
    exploder (IE) or other Windows or Mac OSX specific plugins will not
    work.}. This also supports the scenario of
  thin client, in which the stick is a citrix or VDI client and the
  actual application that is being used is provided by a virtual desktop
  in something like a virtual desktop installation. We know (tested) that
  the Virtual Desktop technology works but also know it is costly in
  license and maintenance. In this mode the stick needs
  \textit{no persistence layer}, and when not provided, that layer
  needs no cleaning, making reuse of the prepared sticks very simple.
\item \textbf{The Unconnected Computer mode.} The stick provides no access except
  to the applications installed. No network, no blue-tooth, no disk,
  nothing. There are two sub use cases:

  \begin{enumerate*}
  \item The application is used as computation aid, the results are
    transferred to paper by the candidate, so no persistent storage
    layer will be required. In such case the stick can be reused without cleaning.
  \item The candidates work is saved in some files and is the actual
    work that needs to be corrected. In this case the work needs to be
    “harvested” from the sticks. We have no experience in this area yet,
    but in theory it is simple with some stick manipulation and some
    simple scripts it should be doable. The same equipment that is used
    to prepare and cleaned can be used for the operation. This is the
    electronic version of handing in paper work by hand.
  \end{enumerate*}
\end{enumerate*}
In the scenarios 1 and 3 the candidates are fully ``sand-boxed'' inside
the OS running on the stick. In the virtual desktop scenario, the
sand-boxing depends on the quality sand-boxing of the virtual desktop.
If the virtual desktop allows e.g. internet access, the stick cannot
prevent it and fraud is just  as possible as using the uncontrolled and
prepared laptop.

\subsubsection*{Current features of the SEBI-Sandbox.}

The following programmes and features are available (status in 2014-03-09).

\begin{itemize*}
\item Boot to graphical desktop without login
\item Network access using functional Fontys account.
\item Developer tools for SEBI  courses: Java SDK and Netbeans or Bluej to
  develop Java programmes.
\item Glassfish JEE container to develop and run java enterprise
  applications.
\item Version control client in graphical and command line versions
\item Two internet browsers: Firefox and google-chrome.
\item Complete postgresql database system with installed exam sample
  database.
\item Command line interface client and pgadmin3 as graphical client to the
  postgres server.
\item Libreoffice calc to be used a spreadsheet/ MS Excel equivalent.
\item Libreoffice base in combination with postgresql as MS Access
  replacement.
\item Xtuple erp client and installed (postgresql) database
\item Widgets to switch keyboard layout for non US keyboard layouts and to
  turn the touch pad on and off during typing.
\end{itemize*}

\subsection*{Special considerations when using examination accounts}
In the \textbf{browser only} or \textbf{kiosk} scenarios,
it is a no-brainer for the candidate to access other information that
is accessible through their normal account on the server. Think of the
N\@tschool scenario, using the normal student account as the candidate
account.This is the case when the normal student account is also used during examinations. 


During examinations, \textbf{\textit{one-time special accounts}} should be used,
which are not to be reused for any other purpose or examination. The
account names and (fresh) password should be provided at the latests
possible moment, to avoid letting Ben en Ahmed share account
information and thus being able to share their information through an
unintended channel, the examination server or upload server. In case of
identical twins, extra care might be warranted.

\subsection*{Restricted access}

Much the same as the USB-stick it selves, the virtual desktop and/or the
special account used should also enforce access  restrictions. This
includes access to local and virtual persistent storage and network.
Examples include the Internet, with services like google internet mail,
drop-box, version control systems like github or even the institute
provided version control system.

\subsubsection*{Naivity is the most enabling factor in fraud}

The most severe form of naivety in this case is ``our students are not
capable of wielding such technology''. Here the speaker forgets that
such ``wielding of technology'' can be purchased from all most anywhere,
by combining things in smart ways.

Venlo, \today
\label{page:last}
\end{document}
